\section{Subset Sums I (Print All Subset Sums)}
\label{sec:rec-subset-sum-1}

\problemstatement{Given an array $\textit{arr}$ of $n$ integers, compute
the \textbf{sum} of every one of its $2^n$ subsets, and return the list
of all such sums (typically in sorted order).}

\begin{patternbox}
A direct simplification of Section~\ref{sec:rec-power-set}'s
pick/not-pick tree: instead of recording the entire subsequence at each
leaf, we only need to record its \textbf{sum} -- a single number threaded
through the recursion as a running total, rather than a list that must
be built and torn down at every level.
\end{patternbox}

\subsection*{Understanding the problem carefully}

Exactly as in Section~\ref{sec:rec-power-set}, every element independently
contributes a pick/not-pick decision. The only change is what we choose
to track: instead of appending the picked element to a growing list (and
later removing it on backtrack), we simply add its value to a running
sum parameter (and, since this is passed by value rather than by
reference, there is nothing to explicitly undo -- each recursive call
gets its own independent copy of the sum).

\begin{ideabox}[Thread a Running Sum Instead of a List]
$\textsc{Solve}(\textit{index}, \textit{sum})$: if
$\textit{index} = n$: record $\textit{sum}$ as one subset's total;
return. Otherwise: $\textsc{Solve}(\textit{index}+1, \textit{sum} +
\textit{arr}[\textit{index}])$ (pick), then
$\textsc{Solve}(\textit{index}+1, \textit{sum})$ (not-pick).
\end{ideabox}

\subsection*{Why no explicit undo step is needed here}

In Section~\ref{sec:rec-power-set}, $\textit{current}$ was a single
shared, mutated list, so the pick branch's append had to be explicitly
undone before exploring the not-pick branch. Here, $\textit{sum}$ is
passed \textbf{by value} (a plain number, copied into each recursive
call rather than referenced and mutated in place), so the pick branch's
modified sum simply ceases to exist once that branch's call returns --
the not-pick branch automatically continues with the original,
unmodified sum from this level, with no explicit ``undo'' arithmetic
required. This is a useful general lesson: whether backtracking needs an
explicit undo step depends entirely on whether the state being threaded
through is mutated in place (needing undo) or passed by value
(automatically restored).

\subsection*{Algorithm}

\begin{enumerate}
  \item $\textsc{Solve}(\textit{index}, \textit{sum})$:
  \begin{enumerate}
    \item If $\textit{index}=n$: append $\textit{sum}$ to the result
          list; return.
    \item $\textsc{Solve}(\textit{index}+1, \textit{sum} +
          \textit{arr}[\textit{index}])$.
    \item $\textsc{Solve}(\textit{index}+1, \textit{sum})$.
  \end{enumerate}
  \item Call $\textsc{Solve}(0,0)$; sort the result list before
        returning (if sorted output is required).
\end{enumerate}

\subsection*{Dry run}

\dryrun{Take $\textit{arr} = [2,3]$.}

\begin{itemize}
  \item $\textsc{Solve}(0,0)$: pick $2$: $\textsc{Solve}(1,2)$.
  \begin{itemize}
    \item Pick $3$: $\textsc{Solve}(2,5)$: base, record $5$.
    \item Not-pick $3$: $\textsc{Solve}(2,2)$: base, record $2$.
  \end{itemize}
  \item Not-pick $2$: $\textsc{Solve}(1,0)$.
  \begin{itemize}
    \item Pick $3$: $\textsc{Solve}(2,3)$: base, record $3$.
    \item Not-pick $3$: $\textsc{Solve}(2,0)$: base, record $0$.
  \end{itemize}
\end{itemize}

Recorded sums: $5,2,3,0$; sorted: $0,2,3,5$ -- matching the four subsets
$\{\}, \{2\}, \{3\}, \{2,3\}$.

\subsection*{C++ implementation}

\begin{lstlisting}
#include <bits/stdc++.h>
using namespace std;

void solve(int index, int sum, vector<int>& arr, vector<int>& result) {
    if (index == (int)arr.size()) {
        result.push_back(sum);
        return;
    }
    solve(index + 1, sum + arr[index], arr, result);
    solve(index + 1, sum, arr, result);
}

vector<int> subsetSums(vector<int>& arr) {
    vector<int> result;
    solve(0, 0, arr, result);
    sort(result.begin(), result.end());
    return result;
}

int main() {
    vector<int> arr = {2, 3};
    for (int s : subsetSums(arr)) cout << s << " ";
    cout << endl;
    return 0;
}
\end{lstlisting}

\begin{complexitybox}
\textbf{Time Complexity.} $O(2^n)$ to generate all sums, plus
$O(2^n \log(2^n)) = O(n \cdot 2^n)$ if sorting the output is required.
\textbf{Space Complexity.} $O(n)$ for the recursion stack, plus
$O(2^n)$ for the output.
\end{complexitybox}

\begin{takeawaybox}
Whether a backtracking branch needs an explicit ``undo'' depends on
whether its state is mutated in place (a shared list, needing an undo)
or passed by value (a plain number or immutable structure, automatically
restored when the branch returns). Recognizing which kind of state a
problem needs is often the deciding factor in how cleanly the recursion
can be written.
\end{takeawaybox}
